\section{Deep Clustering and Contrastive Representation Learning}\label{sec:related:deep_clustering}

Classical clustering assumes that either the raw features or a fixed embedding already reveal the desired partition.
Deep clustering drops that split: the encoder and the discrete structure are learned together.

\Gls{dec} places trainable centroids in embedding space and alternates soft assignment with a Kullback--Leibler objective that sharpens confident assignments toward an auxiliary target \parencite{xie2016unsupervised}.
\Gls{idec} adds a reconstruction term so that the clustering loss does not destroy local geometry \parencite{guo2017idec}.
DeepCluster instead runs \(k\)-means on the current features and treats the assignments as pseudo-labels for a discriminative update \parencite{caron2018deepcluster}.
In all three, the number of clusters is a design choice, and the only partition being learned is a single flat grouping.

Contrastive methods replace explicit centroids with pairwise attraction and repulsion.
InfoNCE and SimCLR pull two views of the same instance together and push other batch members apart \parencite{oord2018representation,chen2020simclr}.
SwAV predicts a small codebook of prototypes from one view and uses the assignment of the other view as the target, which avoids comparing every pair while still preventing collapse \parencite{caron2020swav}.
Contrastive clustering applies the same idea in both the instance dimension and the cluster dimension of a feature matrix \parencite{li2021contrastive}.
These objectives are unsupervised: positives are defined by augmentation, not by labels.

When labels are available, the positive set can be defined by class identity rather than by instance identity.
Supervised contrastive learning does exactly that: every other sample with the same class is a positive, and the loss is an InfoNCE average over those positives \parencite{khosla2020supcon}.
It recovers the triplet loss of FaceNet as a special case that uses a single positive and a hard negative \parencite{schroff2015facenet}.
The method of \cref{sec:method:loss} uses this kernel twice, with different rules for who counts as a positive.

Two properties of this literature matter for the method.
First, training and assignment are usually coupled: \gls{dec}, DeepCluster, and SwAV produce cluster indices as part of the loss.
This thesis decouples them.
The contrastive terms shape the embedding; \gls{dbscan} is applied only at inference and contributes no gradient (\cref{eq:method:cluster_decoding}).
Second, almost every method learns one partition.
A single contrastive term cannot encode both a recording-local emitter grouping and a globally comparable mode grouping, because those two relations contradict each other in one embedding: two pulses from different emitters can share a mode, and two pulses from one emitter can differ in mode.
That conflict is why the encoder splits into two branches after a shared trunk.
